\section*{A costed double auction: mechanism and test}

Q's trading mechanism selects one untraded agent per iteration for at most
\(300\) iterations. A crossing order executes immediately; an improving order
can wait in the book. Q reports that many agents make one-cent improvements,
and that GPT Large's reluctance to cross suppresses trade volume and gross
allocative efficiency (Q, Trading Mechanism; Individual Behavior; Aggressive
Rent-Seeking Suppresses Trade Volume). P explicitly prices generated tokens by
\(C=kTS^{\alpha}\), including reasoning tokens (P, The Energy Society; Small
vs. large models). P shows that removing its size multiplier changes survival,
while larger agents still spend more tokens (P, No Size Penalty). This supplies
an intervention to test in Q, not evidence of its effect in a market.

Let a selected buyer have value \(v\), and suppose the standing ask is
\(a\leq v\). Crossing now gives \(v-a\) before the cost of the current call.
If the buyer instead posts a noncrossing improvement and uses continuation
policy \(\sigma\), let \(F\) indicate a later fill, \(p\) its price, and
\(K_{\sigma}\) the buyer's subsequent inference cost. Let
\(c_{\mathrm{cross}}\) and \(c_{\mathrm{improve}}\) be expected costs of the
current call under the two actions. Crossing is preferred when
\[
v-a-c_{\mathrm{cross}} \geq
\mathbb{E}_{\sigma}\!\left[F(v-p)-K_{\sigma}\right]
-c_{\mathrm{improve}}.
\]
For a flat call fee, the current-call terms cancel; only the expected cost of
continuing differs. Under token metering, current-call token counts may also
differ, so they must be measured. A larger anticipated \(K_{\sigma}\) can
favor crossing, but high cost can also favor abstention or exit. Q's interface
permits declining to announce, yet the trader remains eligible for later calls.
An irrevocable exit option would need to be added in both treatment and control
to make future computation avoidable. The experiment must record declines and
exits separately. This inequality is a mechanism hypothesis, not a behavioral
theorem about language models.

For each executed buyer-seller pair \((b,s)\), the monetary transaction price
cancels from combined trading surplus:
\[
W_{\mathrm{gross}}=\sum_{(b,s)}(v_b-v_s).
\]
For Q's symmetric reservation values the maximum is \(W^{\star}=7.5\):
the six highest buyer values less the six lowest seller costs yield marginal
surpluses \(2.5,2,1.5,1,0.5,0\). Five positive-gain trades can already attain
full gross surplus; Q's sixth equilibrium trade has zero gross gain. With
positive inference cost, inducing that sixth trade can reduce net welfare.
To account for computation in a single-model Q treatment, let \(Y_{i,t}\)
be generated output and reasoning tokens on decision call \(t\), including
calls on which the agent posts no order. Assign a declared conversion
\(\lambda\) from generated tokens to Q's price units and define an
operational net-value measure
\[
W_{\mathrm{net}}(\lambda)=W_{\mathrm{gross}}
-\lambda\sum_{i,t}Y_{i,t}.
\]
P's cost \(kTS^{\alpha}\) uses known local model sizes
\(S\in\{4,8,9\}\) (P, Small vs. large models). Q instead tests
proprietary GPT and Gemini models without disclosing parameter sizes
(Q, Experimental Conditions). The size multiplier cannot therefore be
populated literally for a mixed Q market. Within a fixed-model arm it is
constant and may be absorbed into \(\lambda\). Sweep \(\lambda\) over a
prespecified range that includes costs below, near, and above a price tick
per typical call. Use the same \(\lambda\) to shadow-price the control's
tokens. For cross-model comparison, use a declared model-specific metered
cost schedule and report generated and input tokens separately, because P
counts only generated tokens. P's numerical \(k\) is denominated in
synthetic energy and cannot be copied into Q's prices. This metric is social
welfare only if the conversion reflects real resource opportunity cost; a
fee paid to an exchange is a transfer.

The narrow experiment uses Q's same model, values, order book, horizon, and
matched exogenous selection draws, with \(\lambda=0\) versus explicit metered cost in each
agent's objective. Charge every call, including a no-post call. Add an
irrevocable exit action in both arms, so an agent can avoid further calls. A
flat per-call cost arm can distinguish sensitivity to future calls from
sensitivity to variable token length. Record quote increments, crossing
conditional on a profitable standing quote, declines to announce, permanent
exits, total calls, tokens per completed trade, price dispersion, gross
surplus, and operational net value. Repeat at more than one iteration cap: a
cost effect mediated by waiting should depend on the available horizon. The
decisive result would be an intermediate cost range with earlier crossing and
higher gross surplus whose reduced total token use also improves net value.
A null result or a drop in participation would be informative about the limits
of this mechanism.

The released Q code checks the tick cap, too few active agents, and a
deadlock predicate, but treats an LLM agent as live for that predicate
(\href{https://github.com/jswistak/competitive-market-simulation/blob/master/src/%
market_simulation/graph/nodes/control.py}{released Q implementation}).
Consequently, an early crossing can simply move later calls to remaining
traders; the total may still hit the cap. The proposed exit and stop rule is
therefore a distinct institution and belongs in both experimental arms.
The repository README says its paper configs did not seed mechanism selection;
matched draws require an implementation change
(\href{https://github.com/jswistak/competitive-market-simulation}{repository README}).

The main competing idea is to add P's nonbinding discussion phase to Q's
auction. It is less clean as a first test. P's no-discussion arm has more raw
job collisions, but it also has less deactivation and thus more opportunities
for active agents to choose jobs (P, No Discussion Phase; Discussion). Q's
public order book already transmits offers, and cheap talk might change
information disclosure or strategic manipulation, not just coordination.

Costly offers in human double auctions have already been studied; Jamison and
Plott report slower equilibration and lower efficiency under a tax on bids
and asks
(\href{https://authors.library.caltech.edu/records/fnhzb-n0e58}{Caltech record}).
The publishable comparison must be about LLM inference cost on every call,
the observed one-tick quote loop, and the market's stopping institution,
without presuming a beneficial cost effect.

Limitations: Q and P use different models, objectives, currency units, and
interaction structures. Neither paper establishes that a positive cost will
cause earlier crossing or improve welfare. Q's \(300\)-iteration limit makes
the finite horizon part of the outcome; findings must be checked against a
longer cap. P's five-seed results and model heterogeneity do not justify a
cross-paper model-size claim.
